\chapter[Can Ritual Sustain Unforced Order?]{Can Ritual Propriety Keep Society Running without Coercion?}
\section{A Seating Plan}
Pick up a handwritten seating plan.

For Grandmother's eightieth birthday, the family reserves a restaurant room with a sixteen-person round table. The previous evening Mother repeatedly revises seats: honoured place facing the door, companion at the celebrant's left, precedence between uncles, children near the entrance where dishes pass. She mutters further rules: fish heads face elders, toasts begin with the oldest table, glass rims stay below elders' when clinking.

Who made these rules, where are they written, what happens if we sit wrongly, you ask.

Mother cannot answer any. No gun, law, fine, only perhaps seconds of table silence and relatives' later judgement that a child lacks manners. Yet next day everyone complies, even your usually rebellious cousin, unconsciously lowering his glass to Grandfather without explaining why.

Everyone obeys rules never formally announced. Unwritten, they feel as solid as the tabletop. Elders give them an ancient name: \textbf{li, ritual propriety}.

Look deeper: status, face, first chopsticks, obligations to yield are also seated there. One table miniaturises society: order without police, obedience without commands, shared rules nobody declares.

Our question begins here: \textbf{can society operate chiefly through ritual, conventions, dignity, and shame rather than coercion?}

A man twenty-five centuries ago answered yes and staked his life on testing it: Kong Qiu, Confucius.

\section{A Dance of Sixty-Four}
Enter an early fifth-century BCE day, late Spring and Autumn, at the ruling Jisun ministerial family's ancestral courtyard in Qufu, Lu, today's Shandong.

Mats lie ready for sacrifice, musicians in position, bells and drums striking air. Dancers enter in eight rows of eight, sixty-four young men holding feathers and short flutes. Sleeves fly, spectacle delights guests.

One person is furious.

Ritual yi dance rows follow strict ranks: Son of Heaven eight rows, sixty-four; princes six; ministers four; gentlemen two. Jisun is a minister, entitled to four rows and thirty-two people, yet displays the heavenly king's eight.

The Analects' "Eight Rows" records Confucius on the Ji family:

\begin{quote}
Eight rows dancing in his courtyard! If he can bear to do this, what could he not bear to do?
\end{quote}
Meaning: someone capable of this is capable of anything.

Why such fury over thirty-two extra dancers?

Understand his age. Born 551, died 479 BCE, Confucius lived through old order's dissolution without visible replacement. A nominal Zhou king presided over warring princes governed by strength. Sima Qian's autobiographical preface summarises over two Spring and Autumn centuries: thirty-six murdered rulers, fifty-two destroyed states, innumerable displaced princes unable to preserve their realms.

Collapse spread downward. Kings lost princes, princes ministers. Lu's three hereditary ministerial houses, the Three Huan, eclipsed its ruler, Jisun strongest. Ministers lost retainers: Jisun's Yang Hu imprisoned his master and nearly seized Lu. Every tier usurped another, taking more than entitled. Posterity called it collapsed rites and ruined music.

Thus sixty-four men do not merely dance. They publicly declare dead rules and power's entitlement to any rank. Today dancers, tomorrow the Nine Tripods, next a ruler's head: many murdered princes began losing life through such small things.

Confucius himself came from declining nobility, lost his father early, and said low childhood status taught many humble tasks ("Zi Han"). He held minor storekeeping and livestock offices, reaching Lu's senior justice and security ministry in his fifties. Reforms offended the Three Huan and forced departure. Fourteen years of travelling with pupils persuaded no ruler truly to adopt his programme. He returned to teach late in life, reportedly with three thousand disciples.

Remember this career. Our philosophy is no victor's lesson, but a lifelong political failure's prescription for a failed age.

\section{Can Order Exist without the Whip?}
Now state the question without answering.

Confucius saw one root of disorder: people failing their proper positions and responsibilities. How restore order?

Two visible methods already operated. \textbf{Force}: armies and punishments destroy resistance, while worsening wars produce more chaos. \textbf{Deception}: manipulation, alliances, broken faith, today's pact tomorrow's betrayal. He watched both fail.

His seemingly naive third method: \textbf{install order inside people's bodies}. Replace outer whip with inner ruler, fearful refusal with unwillingness to do wrong. The ruler is li; growing it is learning; its mature bearer is junzi, an exemplary person.

Strong scepticism immediately follows. Would Jisun withdraw dancers because Confucius was angry? No; such powerful families crushed his reforms. Power yields to force, not reasoning. Can taught propriety stop an impolite person's knife?

Europe revisited almost the same problem two millennia later with an opposite answer. Witnessing England's seventeenth-century civil war, Thomas Hobbes's Leviathan finds humanity without overwhelming common power in war of all against all. A strong sovereign maintains peace through fear. Same disease, collapsing order and killing; opposite medicines, a frightening state machine or people needing none.

Who is right? Vote privately before continuing: remove every classroom penalty tomorrow and rely on self-discipline. How long will order last?

\section{Three Voices: Confucius, Mozi, Han Fei}
Pre-Qin thinkers offered at least three serious accounts of order. Hear each fully.

\textbf{First: Confucius, order grown within.}

His programme has several components.

Core is \textbf{ren, benevolence}. An older word becomes central through him. Fan Chi asks its meaning: love people. Zhonggong receives the rule not to impose undesired things on others (both "Yan Yuan"). This is perspective-taking rather than command, using your feelings to measure theirs. Ren is the hidden engine.

Visible \textbf{li}, rites, rules, titles, seats, glass heights, is its operating system: bodily practices install, train, express inner feeling. He fears empty li without ren, asking what ritual can do for an inhumane person ("Eight Rows"). Ren without li hangs unsupported. Asked ritual's root by Lin Fang, he prefers simplicity over luxury and true grief over perfect funerary performance. Reduce form, not sincerity.

He ignites the engine through \textbf{xiao, filial devotion}. Parents are our first others; affection begins its training at home. Disciple You Ruo says those devoted to parents and respectful of brothers rarely offend superiors ("Learning"). Family is ren's first practice ground. Yet xiao is not sentimental ease. Asked by Zixia, Confucius replies se nan, maintaining a pleasant expression is difficult ("Governing"). Money and food are easier than one's face. His disputed injunction against distant travel while parents live, requiring a known destination when travelling ("Li Ren"), returns later.

The formed person is \textbf{junzi}. Confucius quietly transforms lord's son, inherited nobility, into cultivated character open to all and not guaranteed to nobles. A junzi understands rightness, a petty person advantage ("Li Ren"). Birth hierarchy acquires a crack: father no longer determines who you are; cultivation does.

Finally, \textbf{rectification of names}. Asked governance by Duke Jing of Qi, he says ruler should be ruler, minister minister, father father, son son ("Yan Yuan"). Often read as absolute subordination, it actually requires every role, including highest, to \textbf{act like} its name. Rulers failing duties cannot demand ministerial loyalty. Positions form a two-way contract, not a one-way chain.

\textbf{Second: Mozi, who pays for your rituals?}

A generation later, Mozi's "Moderation in Funerals" and "Against Music" attack the programme through \textbf{cost}. Confucian three-year mourning and grand burials bankrupt households, remove labour from production, consume workers' food and clothing through ceremonial music. Moving rituals are luxuries paid by hungry people. Rules that cannot let ordinary people live and eat are decorations, not good rules.

This strikes hard at an issue Confucius less readily confronts: \textbf{ritual has a price, historically paid more heavily by the poor}.

\textbf{Third: Han Fei, both are too optimistic.}

Late Warring States Han Fei overturns the debate. Give him full reasoning, not the villain's mask often imposed on Legalists.

Three layers. Human nature: calculation, not affection, usually motivates. "The Five Vermin" describes a wayward child unmoved by angry parents, rebuking neighbours, instructing teachers, then frightened into reform by armed officials and law. Every soft measure failed; hardness works. Scale: li depends on familiar observers and embarrassment; mobile, large Warring States realms crowd streets with strangers. Who watches? Era: "The Five Vermin" says:

\begin{quote}
High antiquity competed in virtue, the middle age in strategy, the present in strength.
\end{quote}
Han Fei need not love cruelty; selling moral competition in an age of force resembles bringing a spoon to battle.

Three stakes: people can improve; ordinary people must afford rules; most will not improve spontaneously, at least now. Each owns a real problem, none completely refutes another. Their argument continues in new venues.

\section[Thinking Bridge: A Rule's Power]{A Thinking Bridge: Where Does a Rule's Force Come From?}
Take a tool to school, household, legal, group, table rules. Ask three questions:

\textbf{What does it require? What if I refuse, and who imposes the cost? What do I feel when complying?}

The latter two reveal three sources of force:

\begin{itemize}
\item \textbf{Punishment}: lost marks, fines, detention, group expulsion. Fear produces compliance, Han Fei's route.
\item \textbf{Others' eyes}: lost face, embarrassment, damaged image. Shame drives li's main engine.
\item \textbf{Conviction}: no cost question, because you want to act: unseen litter collection, unsupervised honesty. Endorsed values point toward ren.
\end{itemize}
Map your life: traffic rules chiefly punish, seats invite gazes, friendship honesty reflects conviction. Confucius aims not to abolish penalties but move rules inward, from fear to embarrassment to endorsement. \textbf{Learning} moves them.

How is morality learned? The Analects opens with method, not moralising: learning and repeatedly practising bring pleasure ("Learning"). Traditional \mbox{習}, xi, "practise", contains \mbox{羽}, yu, "feathers", above; its original sense is young birds repeatedly learning flight. Practise knowledge until bodily instinct forms. Similar natures diverge through practice, he says ("Yang Huo"). Morality resembles cycling or swimming, learned by doing, not lectures. People are textbooks too: among three walking together, one can find a teacher ("Shu Er"). Exemplars stand everywhere; junzi are living lessons.

This practice-and-example method later divided Confucians who both claimed authentic succession.

Mencius, in the middle Warring States, finds innate moral sprouts. "Gongsun Chou, Part One" offers a famous experiment:

\begin{quote}
Anyone suddenly seeing a child about to fall into a well feels alarm and compassion.
\end{quote}
Not parental favour or fear of blame, but instinctive inability to bear it. Compassion is ren's beginning. Learning therefore \textbf{nurtures} existing sprouts.

Late Warring States Xunzi answers: human nature is bad; goodness is wei ("Human Nature Is Bad"). Wei here means human artifice, not hypocrisy. Desire follows nature; goodness is acquired construction. Sages \textbf{design} li to constrain and transform it. Yet both agree on environment: mugwort among hemp grows straight without support ("Encouraging Learning"). Habit, surroundings, models shape people.

Their civil war concerns \textbf{starting points}, sprout or raw material, but agrees on \textbf{method}: repeated practice and good environments form people. Modern nature-versus-nurture, genes-versus-environment debates still echo it in psychology laboratories.

\section{Read Confucius's Own Answer}
His most direct response appears in the Analects' "Governing":

\begin{quote}
Guide them with ordinances and align them through punishments: people evade penalties without shame. Guide them through virtue and align them through ritual: they gain shame and become correct.
\end{quote}
Dao means guide, qi align. Government commands and penalties encourage evasion without conscience; virtue and li cultivate shame and voluntary correction.

Notice what he \textbf{does not} say: penalties are useless. A former justice minister knew their capacities. More precisely, punishment and li produce different people. Evasion studies loopholes, stopping at red lights for cameras rather than conviction. Internalised standards stop at unobserved junctions. External sensors versus internal scales.

Add his answer to Yan Yuan on ren:

\begin{quote}
Restrain yourself and return to ritual: this is benevolence. For one day do so, and the world returns to benevolence.
\end{quote}
Restrain yourself, not others. Li first constrains its practitioner, a crucial point repeatedly forgotten later.

Honestly, these are \textbf{proposals}, not established facts. Confucius had no survey data showing ritual guidance reduced crime more than punishment; we cannot invent it for him. He bets on human mouldability. Later accounts weigh that wager.

\section{One Li, Three Explanations}
Establish a broadly accepted fact: traditional Chinese local disputes mostly stayed outside government. Inheritance, loans, boundaries, marriage were settled by clans, village compacts, mediators. Litigation could mean shame or ruin; absence of lawsuits signalled good rule. Ming--Qing gazetteers, genealogies, contracts document this mechanism.

One fact, at least three unequal explanations.

\textbf{Explanation 1: Ritual governance is low-cost administration.}

Officials reached county level, with one magistrate over hundreds of thousands and no affordable coercive network for every village. Li fills the gap through clan heads, compacts, face, self-government reaching tables without police or courts. This explains dynasties' enthusiasm for moral commendations and filial or chaste memorial arches: rulers like cheap order. But cheap for \textbf{whom}? Lower governance costs went somewhere; this account does not ask where.

\textbf{Explanation 2: Ritual morality is refined oppression.}

The weaker paid. In Lu Xun's 1918 "Diary of a Madman", the protagonist reads history:

\begin{quote}
I opened history and found no dates. On every crooked page stood benevolence, righteousness, and morality. Sleepless, I studied half the night, until between the lines I found two words filling the book: eat people!
\end{quote}
Here eating people means moral language wrapping oppression as virtue until even victims' anger appears immature. Ming--Qing chastity arches glorified widows who never remarried, funded by clans and praised by courts. No officer whipped them into widowhood. Precisely: \textbf{li needs no whip when a twenty-year-old widow enters her own cage and society, perhaps herself included, willingly guards it}. From another position, no-litigation ideals tell victims to endure silently. This sees hidden costs, but cannot explain genuine ordinary needs if every ritual is oppression: funerals settle real grief; ruler-as-ruler expectations constrain some rulers. All-power interpretations miss emotion.

\textbf{Explanation 3: Li operates familiar society.}

Take structural rather than court or victim perspective. Lifelong neighbours in immobile villages own reputation above all; shame is practical punishment. Li is neither sages' gift nor rulers' invented conspiracy but long association's operating system. Everyone enforces violations by withholding loans, marriage ties, advocacy. This explains cross-cultural equivalents. Louis XIV enclosed nobles at Versailles through elaborate etiquette, shirt-presenting privileges, stools versus standing, transforming territorial lords into favour-seeking courtiers. France independently invented ritual rule, indicating structure rather than Chinese national character. Yet spontaneous-growth accounts underplay design and reform, and struggle with surviving seating plans after dispersed families cease being close neighbours.

Each explanation holds one ledger page. Purely good or bad li usually means other pages closed prematurely.

\section[Further Challenge: Ritual Changes People]{A Deeper Challenge: Why Can Ritual Change People?}
Place Confucius's intuition under a modern sociological lens associated with Émile Durkheim.

Studying Aboriginal Australian rites early in the twentieth century, he noted normally dispersed people gathering to dance, chant, dress alike, and enter collective excitement. Doing \textbf{the same action} generates we, the moment society itself forms. Rites do not merely \textbf{express} prior unity; they \textbf{produce} it. Without periodic shared action, belonging disperses.

Reconsider school anthems, stadium waves, precisely folded military bedding and marching, kin-ranked funerals, classroom greetings. Why do states, armies, teams, schools, families repeat useless acts? Durkheim calls them organisational chargers. A class never doing anything together lacks an actual we.

Confucius now sounds modern. Joining rites and music, insisting on sacrifices, coming-of-age ceremonies, funerals, seemed conservative. Sociologically, he found the manufacture of social glue, perhaps understanding dancers' danger better than later critics. Jisun was \textbf{privately producing} we, charging his family through royal ritual with energy stolen from Lu's ruling house.

The theory also exposes risk: creating us delineates them. New pupils cannot yet sing; some relatives never reach honoured seats. Warm interiors harden boundaries. Japanese bowing angles, honorific levels, after-work drinks strongly bind society while marginalising foreigners, transfer pupils, socially awkward people. Ritual's cohesion and exclusion outlets run simultaneously.

Where does Durkheim stand in the Confucian dispute? Apparently with Xunzi: society installs much of the self. Mencius replies that the child's falling-well shock arrives before ritual, suggesting something not installed. Social construction or awakenable nature? A century after Durkheim, renamed debate continues.

\section{Class Agreements and Family Chats}
This ancient question lives in classrooms and WeChat.

Three disciplinary devices operate: school penalties, agreed class duties and library rules, and powerful invisible class ethos, we do not do that. Some teachers first construct belonging through won matches, slogans, shared memories, making violations embarrassing. This small Confucian experiment sometimes works: you have probably seen orderly unsupervised study.

It has a dark version. We-do-not may suppress not misconduct but defenceless newcomers, poor children unable to join activities, ungregarious personalities. Everyone enforces the gaze; victims cannot identify whom to report. Nobody strikes, everyone naturally stops speaking. Lu Xun's cage exists at class scale.

Family holds the sharpest tension. Separate \textbf{Confucius's words, later transformations, and today's evaluation}.

Parents alive, avoid distant travel, or state your destination. Yet fourteen-year-olds may soon study elsewhere; parents may already work elsewhere. Can no-distance survive mass mobility? What then of New Year travel? Later filial teaching includes Yuan-compiled Twenty-Four Filial Exemplars' Guo Ju, poor enough to plan burying his son alive to preserve his mother's food. Taught for centuries, it made young Lu Xun afraid of becoming filial. Compare pleasant expression with killing one's child for parents: ethics can grow crooked from within.

Give Confucius fairness too: he did not make filiality absolute obedience. "Li Ren" says:

\begin{quote}
Serve parents with gentle remonstrance. If they do not follow your intention, remain respectful without defiance, troubled but not resentful.
\end{quote}
Meaning: tactfully advise wrongdoing; if unheard, preserve respect without confrontation or resentment. Its structure makes \textbf{remonstrance part of filiality, not obedience}. Disagreement may be a duty; avoid humiliation in expressing it. Closer to modern equality than imagined, its deeper difference concerns hierarchy: modern persons hold equal standing regardless of age; Confucian relations naturally slope by seniority, equality not sameness. Adults-speak-children-silent versus I-am-a-family-member debates that line.

Anonymous online comments finally stress-test everything: weak punishment through replaceable accounts, weak gaze without identities, undeveloped conviction. Han Fei points to predictable results with two engines removed. Confucius might ask whether decent public behaviour collapsing in anonymity reveals installed moral rulers or merely surveillance sensors.

\section[Your Turn: A Class without Penalties]{Your Turn: The Class That Cancels Every Punishment}
Return to your earlier vote.

Next term your class pilots no lost marks, standing punishments, parent summons. Only shared agreements, recognised exemplars, and genuine embarrassment at harming the group remain. Confucius in a classroom: would you bet on success?

Lay every counter out first.

Confucius has observed self-discipline, exemplary groups, embarrassment outlasting fear. Habit and environment genuinely mould people, agreed by Mencius and Xunzi despite their dispute.

Han Fei has the one or two shameless classmates unaffected by soft constraints, capable of destroying everyone's normative gaze; one broken window lets rules leak.

Mozi asks who can afford the experiment. Do agreements protect excluded children or majority comfort?

Lu Xun weighs heaviest: shame is power, and power presses easiest targets. Sensitive, conscientious children pay highest compliance costs while indifferent peers escape. Shame-based punishment may be exceptionally unfairly distributed.

Now answer: can li sustain society without coercion?

If yes, what size, familiarity, cultivation does it need, and can today's world supply them? If no, what cost for more cameras, finer penalty tables, denser laws? Would you choose self-discipline defined by surrounding gazes, or rules indifferent to inner belief?

No standard answer. Confucius staked his life on one side; Han Fei on another, dying in the Qin prison of the state he served. Many bet, including you. Give your answer and reasons.
